\section{CEO 与 Board：建立不失真的信息回路}

公司扩大后，CEO 不再是所有问题的直接解决者，而是信息、资本、人才和风险的汇聚点。角色的“孤独”常不是没人说话，而是坏消息在层级中被过滤、团队期待确定答案、董事会又只能基于管理层提供的 evidence 做判断。系统解法是明确 decision rights 和 unfiltered information loop。

\subsection{Bad news、options 与 accountability}

管理团队应把 operating truth 转成可讨论对象：指标发生了什么，客户或安全风险在哪里，哪些 root cause 已证实，哪些只是 hypothesis，有哪些 options、trade-off 和 time constraint。董事会的作用不是替管理层运行公司，而是挑战假设、配置资本和监督重大风险；CEO 则对最终选择与 follow-through 负责。

\lecturefigure{07-ceo-board-loop.png}{CEO 与董事会需要从 operating truth 到 options、decision 和 follow-through 的无失真回路。}{本地字幕 12:00--14:40；概念重绘。}

\begin{knowledgebox}{读图：Dashboard 不能代替坏消息叙事}
Metrics 告诉董事会“哪里异常”，customer evidence 与 incident timeline 解释“为什么重要”，options 明确“可以做什么”，owner/milestone 让决策可复查。如果 CEO 只报告已经解决的问题，board 会在真正需要支持时才第一次接触上下文；如果 board 越过 CEO 直接管理执行，又会破坏责任链。
\end{knowledgebox}

\teachervoice{讲者回顾早期董事会互动时，认为自己曾过度想证明“一切都在掌控中”。后来更有效的方法是把困难和坏消息提前带入讨论。对 frontier systems 团队而言，这等价于让风险在还能选择时出现，而不是在 outage 或 breach 后才升级。}

\subsection{本章小结}

CEO/board loop 的质量取决于信息完整、option clarity 与责任边界。孤独不能靠更多 status meeting 解决，而要靠能容纳不确定性和坏消息的治理结构。

